% ═══════════════════════════════════════════════════════════════════════════
%  فصل ۲ — مبانی حقوقی مداخله‌ی خارجی در حقوق بین‌الملل
% ═══════════════════════════════════════════════════════════════════════════
\chapter{مبانی حقوقی مداخله‌ی خارجی در حقوق بین‌الملل}
\label{ch:legal}

\begin{tcolorbox}[mybox, title={چکیده‌ی فصل}]
این فصل به بررسی نظام حقوقی بین‌المللی حاکم بر مداخله‌ی خارجی می‌پردازد. از اصل حاکمیت وستفالیایی تا دکترین «مسئولیت حمایت» (\lr{R2P})، تحول هنجارهای بین‌المللی، رویه‌ی دیوان بین‌المللی دادگستری، و مشروعیت حقوقی «مداخله بر اساس دعوت» بررسی می‌شود.
\end{tcolorbox}

% ─────────────────────────────────────────────────────────────────────────
\section{اصل حاکمیت دولت‌ها و منع مداخله}
\label{sec:ch02-sovereignty}


\subsection{صلح وستفالی و تولد اصل حاکمیت}
\label{subsec:ch02-westphalia}

نظام مدرن دولت‌ـ‌ملت‌ها ریشه در \histref{صلح وستفالی}{۱۶۴۸ م.} دارد. این معاهده که به جنگ‌های سی‌ساله در اروپا پایان داد، دو اصل بنیادین را بنا نهاد:

\begin{enumerate}[nosep, label=\textbf{\arabic*.}]
  \item \concept{اصل برابری حاکمیت}: هر دولت در قلمرو خود اقتدار عالی دارد و هیچ قدرت فراملی حق دخالت در امور داخلی آن را ندارد.
  \item \concept{اصل عدم مداخله}: دولت‌ها حق ندارند در امور داخلی یکدیگر مداخله کنند.%
  \footnote{\lr{Krasner, Stephen D. \textit{Sovereignty: Organized Hypocrisy}. Princeton UP, 1999, pp.~20--25.}}
\end{enumerate}

\thinker{ژان بُدَن} (\lr{Jean Bodin}) پیش از وستفالی، حاکمیت را «قدرت مطلق و دائمی یک جمهوری» تعریف کرده بود.%
\footnote{\lr{Bodin, Jean. \textit{Six Books of the Commonwealth}. Trans. M.J. Tooley. Blackwell, 1955 [1576], Book I, Ch.~8.}}
اما در عمل، اصل حاکمیت هرگز مطلق نبوده و از همان آغاز با استثنائاتی همراه بوده است. \thinker{هوگو گروسیوس} (\lr{Hugo Grotius})، پدر حقوق بین‌الملل، حق مداخله‌ی بشردوستانه را در صورت ستم فاحش حاکم بر رعایایش به رسمیت می‌شناخت.%
\footnote{\lr{Grotius, Hugo. \textit{De Jure Belli ac Pacis}. Trans. F.W. Kelsey. Clarendon, 1925 [1625], Book II, Ch.~25.}}

\begin{legalNote}{اصل حاکمیت در منشور سازمان ملل}
ماده‌ی ۲ بند ۱ منشور سازمان ملل متحد (۱۹۴۵) مقرر می‌دارد: «سازمان بر مبنای اصل برابری حاکمیت کلیه‌ی اعضا استوار است.» بند ۷ همین ماده تصریح می‌کند: «هیچ‌یک از مقررات منشور حاضر، ملل متحد را مجاز نمی‌دارد در اموری که اساساً در صلاحیت داخلی هر دولتی است مداخله نمایند.»%
\footnote{منشور سازمان ملل متحد، ماده‌ی ۲، بندهای ۱ و ۷. متن رسمی فارسی.}

با این حال، فصل هفتم منشور (مواد ۳۹ تا ۵۱) به شورای امنیت اجازه می‌دهد در صورت «تهدید علیه صلح، نقض صلح یا عمل تجاوزکارانه» اقدامات اجباری — از جمله نظامی — اتخاذ کند. این استثنا، روزنه‌ی اصلی برای مداخلات بین‌المللی «مشروع» بوده است.
\end{legalNote}

\subsection{تحول تاریخی اصل عدم مداخله}
\label{subsec:ch02-evolution}

اصل عدم مداخله در طول تاریخ حقوق بین‌الملل تحولات مهمی را پشت سر گذاشته است:

\begin{figure}[htbp]
\centering
\begin{tikzpicture}[
  x=0.016\linewidth,
  every node/.style={font=\scriptsize},
]

% خط اصلی
\draw[line width=2pt, PrimaryMid] (0,0) -- (130,0);

% نقاط زمانی
\node[circle, fill=EraAncient, minimum size=8pt, inner sep=0pt] (w) at (0,0) {};
\node[circle, fill=EraMedieval, minimum size=8pt, inner sep=0pt] (v) at (25,0) {};
\node[circle, fill=EraEarlyMod, minimum size=8pt, inner sep=0pt] (h) at (45,0) {};
\node[circle, fill=EraModern, minimum size=8pt, inner sep=0pt] (u) at (65,0) {};
\node[circle, fill=EraContemp, minimum size=8pt, inner sep=0pt] (n) at (85,0) {};
\node[circle, fill=EraPostmod, minimum size=8pt, inner sep=0pt] (r) at (105,0) {};
\node[circle, fill=AccentRed, minimum size=8pt, inner sep=0pt] (s) at (125,0) {};

% برچسب‌ها بالا
\node[above=0.6cm, text width=2.2cm, align=center] at (w) {\rl{وستفالی}\\1648};
\node[above=0.6cm, text width=2.5cm, align=center] at (v) {\rl{کنگره‌ی وین}\\1815};
\node[above=0.6cm, text width=2.5cm, align=center] at (h) {\rl{کنوانسیون‌های لاهه}\\1899/1907};
\node[above=0.6cm, text width=2.5cm, align=center] at (u) {\rl{منشور ملل متحد}\\1945};
\node[above=0.6cm, text width=2.5cm, align=center] at (n) {\rl{قطعنامه ۲۶۲۵}\\1970};
\node[above=0.6cm, text width=2cm, align=center] at (r) {\rl{اجلاس جهانی}\\2005};
\node[above=0.6cm, text width=2cm, align=center] at (s) {\rl{بحران لیبی}\\2011};

% برچسب‌ها پایین
\node[below=0.6cm, text width=2.8cm, align=center, text=NeutralDark] at (w)
  {\rl{تولد اصل حاکمیت}};
\node[below=0.6cm, text width=2.8cm, align=center, text=NeutralDark] at (v)
  {\rl{مداخله‌ی جمعی مشروع}};
\node[below=0.6cm, text width=2.8cm, align=center, text=NeutralDark] at (h)
  {\rl{حقوق جنگ}};
\node[below=0.6cm, text width=2.8cm, align=center, text=NeutralDark] at (u)
  {\rl{منع توسل به زور}};
\node[below=0.6cm, text width=2.8cm, align=center, text=NeutralDark] at (n)
  {\rl{حق تعیین سرنوشت}};
\node[below=0.6cm, text width=2.8cm, align=center, text=NeutralDark] at (r)
  {\rl{مسئولیت حمایت}};
\node[below=0.6cm, text width=2.8cm, align=center, text=NeutralDark] at (s)
  {\rl{اِعمال \lr{R2P}}};

\end{tikzpicture}
\caption{سیر تحول هنجاری اصل عدم مداخله و استثنائات آن}
\label{fig:ch02-evolution-timeline}
\end{figure}

% ─────────────────────────────────────────────────────────────────────────
\section{مداخله بر اساس دعوت (\texorpdfstring{\lr{Intervention by Invitation}}{Intervention by Invitation})}
\label{sec:ch02-invitation}
% ─────────────────────────────────────────────────────────────────────────

\subsection{مبانی حقوقی}
\label{subsec:ch02-invitation-legal}

مداخله بر اساس دعوت (\lr{Intervention by Invitation}) یکی از پیچیده‌ترین موضوعات حقوق بین‌الملل است. اصل کلی آن است که اگر \concept{حکومت مشروع} یک دولت، از دولت دیگری برای مداخله — اعم از نظامی یا غیرنظامی — در امور داخلی‌اش دعوت کند، این مداخله از نظر حقوقی مجاز است؛ زیرا \concept{رضایت دولت} استثنایی بر اصل منع توسل به زور (ماده‌ی ۲ بند ۴ منشور) محسوب می‌شود.%
\footnote{\lr{Doswald-Beck, Louise. "The Legal Validity of Military Intervention by Invitation of the Government." \textit{British Year Book of International Law} 56, no.~1 (1985): 189--252.}}

\thinker{لوئیز دوسوالد-بک} (\lr{Doswald-Beck}) در پژوهش بنیادین خود، شرایط مشروعیت حقوقی مداخله بر اساس دعوت را چنین برشمرده است:

\begin{tcolorbox}[legalbox, title={شرایط حقوقی مداخله بر اساس دعوت}]
\begin{enumerate}[nosep, label=\textbf{شرط \arabic*:}]
  \item \textbf{مشروعیت دعوت‌کننده:} دعوت باید از سوی «حکومت مشروع» — یعنی حکومتی که از نظر حقوق بین‌الملل به رسمیت شناخته شده — صادر شود.
  \item \textbf{رضایت واقعی:} رضایت نباید تحت اجبار، فریب یا فشار نظامی مداخله‌گر باشد.
  \item \textbf{تناسب:} مداخله باید متناسب با درخواست و محدود به هدف اعلام‌شده باشد.
  \item \textbf{عدم نقض حقوق بنیادین:} مداخله نباید ناقض هنجارهای آمره (\lr{jus cogens}) — مانند ممنوعیت نسل‌کشی — باشد.
  \item \textbf{عدم مداخله در جنگ داخلی:} اگر کشور هدف درگیر جنگ داخلی تمام‌عیار باشد، مشروعیت دعوت حکومت زیر سؤال می‌رود.
\end{enumerate}
\end{tcolorbox}

\subsection{رویه‌ی دیوان بین‌المللی دادگستری}
\label{subsec:ch02-icj}

دیوان بین‌المللی دادگستری (\lr{ICJ}) در چندین پرونده به مسئله‌ی مداخله بر اساس دعوت پرداخته است. مهم‌ترین آنها عبارت‌اند از:

\begin{tcolorbox}[casebox, title={پرونده‌ی نیکاراگوئه علیه ایالات متحده (۱۹۸۶)}]
در این پرونده، دیوان مقرر داشت که «اصل عدم مداخله، حق هر دولت حاکم را برای اداره‌ی امور خود بدون دخالت خارجی تضمین می‌کند.» دیوان تأکید کرد که حمایت از گروه‌های مسلح اپوزیسیون (کنتراها) علیه حکومت نیکاراگوئه، نقض حقوق بین‌الملل بوده و دفاع امریکا مبنی بر «دعوت» از سوی اپوزیسیون مردود است.%
\footnote{\lr{ICJ. \textit{Military and Paramilitary Activities in and against Nicaragua (Nicaragua v. United States of America)}. Judgment, 27 June 1986. ICJ Reports 1986, p.~14, paras.~202--209.}}

نکته‌ی کلیدی: دیوان صراحتاً اعلام کرد که «دعوت» از سوی اپوزیسیون (نه حکومت مشروع) مبنای حقوقی برای مداخله نظامی ایجاد نمی‌کند.
\end{tcolorbox}

\begin{tcolorbox}[casebox, title={پرونده‌ی فعالیت‌های مسلحانه در قلمرو کنگو (۲۰۰۵)}]
در این پرونده، دیوان رضایت دولت جمهوری دموکراتیک کنگو به حضور نظامی اوگاندا را بررسی کرد و نتیجه گرفت که حتی اگر رضایت اولیه وجود داشته باشد، تداوم حضور نظامی پس از لغو رضایت غیرقانونی است.%
\footnote{\lr{ICJ. \textit{Armed Activities on the Territory of the Congo (Democratic Republic of the Congo v. Uganda)}. Judgment, 19 December 2005. ICJ Reports 2005, p.~168, paras.~42--53.}}
\end{tcolorbox}

\subsection{معضل «حکومت مشروع» در جنگ‌های داخلی}
\label{subsec:ch02-legitimate-gov}

یکی از بزرگ‌ترین چالش‌های حقوقی مداخله بر اساس دعوت، تعیین «حکومت مشروع» در شرایط جنگ داخلی است. \thinker{جیمز کرافورد} (\lr{James Crawford}) و \thinker{مارکو ساسولی} (\lr{Marco Sassòli}) نشان داده‌اند که در عمل، «مشروعیت» اغلب بر اساس \concept{واقعیت سیاسی} (کنترل مؤثر بر قلمرو) و نه بر اساس معیارهای دموکراتیک تعیین می‌شود.%
\footnote{\lr{Crawford, James. \textit{Brownlie's Principles of Public International Law}. 9th ed. Oxford UP, 2019, pp.~142--148.}}

\begin{longtable}{
  >{\raggedleft\arraybackslash}p{2.5cm}
  >{\raggedleft\arraybackslash}p{2.5cm}
  >{\raggedleft\arraybackslash}p{3cm}
  >{\raggedleft\arraybackslash}p{2.5cm}
  >{\raggedleft\arraybackslash}p{3cm}}
\caption{نمونه‌های بحث‌برانگیز مداخله بر اساس دعوت}\label{tab:ch02-cases}\\
\toprule
\rowcolor{PrimaryDeep}
\textcolor{white}{\textbf{مورد}} &
\textcolor{white}{\textbf{دعوت‌کننده}} &
\textcolor{white}{\textbf{مداخله‌گر}} &
\textcolor{white}{\textbf{سال}} &
\textcolor{white}{\textbf{وضعیت حقوقی}} \\
\midrule
\endfirsthead
\toprule
\rowcolor{PrimaryDeep}
\textcolor{white}{\textbf{مورد}} &
\textcolor{white}{\textbf{دعوت‌کننده}} &
\textcolor{white}{\textbf{مداخله‌گر}} &
\textcolor{white}{\textbf{سال}} &
\textcolor{white}{\textbf{وضعیت حقوقی}} \\
\midrule
\endhead
\bottomrule
\endlastfoot

مجارستان &
دولت کادار &
شوروی &
۱۹۵۶ &
بحث‌برانگیز%
\footnote{\lr{Franck, Thomas M. \textit{Recourse to Force}. Cambridge UP, 2002, pp.~95--98.}} \\
\rowcolor{NeutralLight}
چکسلواکی &
«درخواست رفقا» &
شوروی (پیمان ورشو) &
۱۹۶۸ &
بسیار بحث‌برانگیز%
\footnote{دکترین برژنف مبنای ادعای شوروی بود. رجوع کنید به: \lr{Ouimet, Matthew J. \textit{The Rise and Fall of the Brezhnev Doctrine}. U of North Carolina P, 2003.}} \\
افغانستان &
دولت کارمل &
شوروی &
۱۹۷۹ &
محکوم‌شده (قطعنامه مجمع عمومی)%
\footnote{\lr{UN General Assembly Resolution ES-6/2. 14 January 1980.}} \\
\rowcolor{NeutralLight}
گرانادا &
فرماندار کل &
ایالات متحده &
۱۹۸۳ &
بحث‌برانگیز \\
سوریه &
دولت اسد &
روسیه و ایران &
۲۰۱۵ &
از نظر فنی مجاز؛ اخلاقاً بحث‌برانگیز%
\footnote{\lr{Ruys, Tom, and Luca Ferro. "Weathering the Storm: Legality and Legal Implications of the Saudi-Led Military Intervention in Yemen." \textit{ICLQ} 65 (2016): 61--98.}} \\
\rowcolor{NeutralLight}
مالی &
دولت مالی &
فرانسه &
۲۰۱۳ &
نسبتاً پذیرفته‌شده%
\footnote{\lr{Bothe, Michael. "The Legality of the French Intervention in Mali." \textit{EJIL Talk}, January 2013.}} \\

\end{longtable}

% ─────────────────────────────────────────────────────────────────────────
\section{دکترین مسئولیت حمایت (\texorpdfstring{\lr{R2P}}{R2P})}
\label{sec:ch02-r2p}
% ─────────────────────────────────────────────────────────────────────────

\subsection{زمینه‌ی تاریخی: از رواندا تا کوزوو}
\label{subsec:ch02-r2p-background}

شکست جامعه‌ی بین‌المللی در جلوگیری از \enemy{نسل‌کشی رواندا} (\histref{}{۱۹۹۴}) و \enemy{قتل‌عام سربرنیتسا} (\histref{}{۱۹۹۵}) و سپس مداخله‌ی بدون مجوز شورای امنیت ناتو در \histref{کوزوو}{۱۹۹۹}، بحران عمیقی در حقوق بین‌الملل ایجاد کرد. \thinker{کوفی عنان}، دبیرکل وقت سازمان ملل، در گزارش هزاره (۲۰۰۰) پرسش بنیادینی مطرح کرد:

\begin{histQuote}{کوفی عنان، ۲۰۰۰}
«اگر مداخله‌ی بشردوستانه واقعاً تعرض غیرقابل قبولی به حاکمیت است، پس ما باید به مردم رواندا و سربرنیتسا چه پاسخ دهیم — مردمی که حاکمیت دولت‌هایشان ابزار سرکوب‌شان بود؟»%
\footnote{\lr{Annan, Kofi. \textit{We the Peoples: The Role of the United Nations in the 21st Century}. UN, 2000, p.~48.}}
\end{histQuote}

\subsection{گزارش کمیسیون بین‌المللی (۲۰۰۱)}
\label{subsec:ch02-iciss}

در پاسخ به این پرسش، دولت کانادا \concept{کمیسیون بین‌المللی مداخله و حاکمیت دولت} (\lr{ICISS}) را تأسیس کرد. گزارش نهایی این کمیسیون (دسامبر ۲۰۰۱) مفهوم «مسئولیت حمایت» را با سه رکن معرفی کرد:

\begin{tcolorbox}[legalbox, title={سه رکن مسئولیت حمایت}]
\begin{enumerate}[nosep, label=\textbf{رکن \arabic*:}]
  \item \textbf{مسئولیت پیشگیری} (\lr{Responsibility to Prevent}): رسیدگی به ریشه‌های بحران پیش از وقوع فاجعه.
  \item \textbf{مسئولیت واکنش} (\lr{Responsibility to React}): اقدام قاطع — شامل مداخله‌ی نظامی در موارد حاد — هنگامی که دولت ناتوان یا ناخواسته از حمایت شهروندانش باشد.
  \item \textbf{مسئولیت بازسازی} (\lr{Responsibility to Rebuild}): کمک به بازسازی و صلح‌سازی پس از مداخله.
\end{enumerate}
\vspace{6pt}
\textit{منبع:} \lr{ICISS. \textit{The Responsibility to Protect}. IDRC, 2001.}%
\footnote{\lr{International Commission on Intervention and State Sovereignty (ICISS). \textit{The Responsibility to Protect}. IDRC, 2001, pp.~xi--xiii.}}
\end{tcolorbox}

\subsection{سند اجلاس جهانی ۲۰۰۵}
\label{subsec:ch02-world-summit}

در اجلاس جهانی ۲۰۰۵، سران دولت‌ها بندهای ۱۳۸ و ۱۳۹ سند نهایی را تصویب کردند که \lr{R2P} را در چهار جرم محدود ساخت:

\begin{tcolorbox}[enemybox, title={چهار جرم مشمول مسئولیت حمایت}]
\begin{enumerate}[nosep, label=\textbf{\arabic*.}]
  \item \enemy{نسل‌کشی} (\lr{Genocide})
  \item \enemy{جنایات جنگی} (\lr{War Crimes})
  \item \enemy{پاک‌سازی قومی} (\lr{Ethnic Cleansing})
  \item \enemy{جنایات علیه بشریت} (\lr{Crimes against Humanity})
\end{enumerate}
\vspace{4pt}
\textit{بند ۱۳۹:} «ما آماده‌ایم که از طریق شورای امنیت و بر مبنای فصل هفتم منشور، اقدام جمعی به‌موقع و قاطعانه‌ای انجام دهیم... هرگاه ابزارهای مسالمت‌آمیز ناکافی باشد و مقامات ملی آشکارا از حمایت مردم خود ناتوان باشند.»%
\footnote{\lr{UN General Assembly. \textit{2005 World Summit Outcome}. A/RES/60/1, paras.~138--139.}}
\end{tcolorbox}

\subsection{نقد \texorpdfstring{\lr{R2P}}{R2P}: مشروعیت‌بخشی به امپریالیسم؟}
\label{subsec:ch02-r2p-critique}

دکترین \lr{R2P} از سوی بسیاری از کشورهای جهان سوم و اندیشمندان پسااستعماری مورد نقد قرار گرفته است:

\begin{itemize}[nosep, rightmargin=1em]
  \item \thinker{نوام چامسکی} (\lr{Noam Chomsky}) آن را «نسخه‌ی نوین بار سفیدپوست» (\lr{White Man's Burden}) می‌خواند.%
  \footnote{\lr{Chomsky, Noam. \textit{A New Generation Draws the Line}. Verso, 2012.}}
  \item \thinker{محمود ممدانی} (\lr{Mahmood Mamdani}) استدلال می‌کند که \lr{R2P} ابزاری برای «تقسیم جهان به ناجیان و قربانیان» است.%
  \footnote{\lr{Mamdani, Mahmood. "Responsibility to Protect or Right to Punish?" \textit{Journal of Intervention and Statebuilding} 4, no.~1 (2010): 53--67.}}
  \item نمایندگان چین و روسیه در شورای امنیت بارها \lr{R2P} را «ابزار تغییر رژیم» خوانده‌اند — به‌ویژه پس از تجربه‌ی لیبی (۲۰۱۱).%
  \footnote{\lr{Zifcak, Spencer. "The Responsibility to Protect after Libya and Syria." \textit{Melbourne Journal of International Law} 13, no.~1 (2012): 1--35.}}
\end{itemize}

% ─────────────────────────────────────────────────────────────────────────
\section{حق تعیین سرنوشت و مداخله‌ی خارجی}
\label{sec:ch02-self-determination}
% ─────────────────────────────────────────────────────────────────────────

\subsection{ریشه‌ها و تحول مفهوم}
\label{subsec:ch02-sd-evolution}

حق تعیین سرنوشت (\lr{Self-Determination}) از اصول بنیادین حقوق بین‌الملل معاصر است. ماده‌ی ۱ مشترک میثاقین بین‌المللی (۱۹۶۶) مقرر می‌دارد: «همه‌ی مردمان حق تعیین سرنوشت دارند. به موجب این حق، آنان آزادانه وضعیت سیاسی خود را تعیین و آزادانه تحول اقتصادی، اجتماعی و فرهنگی خود را دنبال می‌کنند.»%
\footnote{میثاق بین‌المللی حقوق مدنی و سیاسی، ماده‌ی ۱.}

\begin{table}[htbp]
\centering
\caption{تحول مفهوم حق تعیین سرنوشت}\label{tab:ch02-self-det}
\begin{adjustbox}{max width=\linewidth}
\begin{tabular}{
  >{\raggedleft\arraybackslash\bfseries}p{2.5cm}
  >{\raggedleft\arraybackslash}p{3cm}
  >{\raggedleft\arraybackslash}p{4cm}
  >{\raggedleft\arraybackslash}p{4cm}}
\toprule
\rowcolor{PrimaryDeep}
\textcolor{white}{دوره} &
\textcolor{white}{مبنا} &
\textcolor{white}{دامنه} &
\textcolor{white}{محدودیت} \\
\midrule
۱۹۱۸–۱۹۴۵ &
چهارده اصل ویلسون &
ملت‌های اروپایی &
عملاً محدود به متفقین \\
\rowcolor{NeutralLight}
۱۹۴۵–۱۹۶۰ &
منشور ملل متحد (ماده ۱ بند ۲) &
ملت‌های تحت استعمار &
استعمارزدایی \\
۱۹۶۰–۱۹۷۰ &
قطعنامه ۱۵۱۴ و ۲۶۲۵ &
همه‌ی مردمان تحت سلطه‌ی بیگانه &
منع تجزیه‌ی ارضی دولت‌ها \\
\rowcolor{NeutralLight}
۱۹۹۰–کنون &
رویه‌ی دیوان و عملکرد دولت‌ها &
حق تعیین سرنوشت داخلی (خودمختاری) &
جدایی تنها در «شرایط استثنایی»%
\footnote{\lr{Supreme Court of Canada. \textit{Reference Re Secession of Quebec}. [1998] 2 SCR 217, para.~134.}} \\
\bottomrule
\end{tabular}
\end{adjustbox}
\end{table}

\subsection{تنش میان حق تعیین سرنوشت و تمامیت ارضی}
\label{subsec:ch02-tension}

تنش بنیادین حقوقی آنجاست که حق تعیین سرنوشت ممکن است با اصل تمامیت ارضی تعارض پیدا کند. این تنش در موارد زیر بارز بوده است:

\begin{itemize}[nosep, rightmargin=1em]
  \item \textbf{کوزوو (۲۰۰۸):} اعلام استقلال یک‌جانبه؛ نظر مشورتی دیوان (۲۰۱۰) آن را «مغایر با حقوق بین‌الملل» ندانست اما حق جدایی یک‌جانبه را نیز تأیید نکرد.%
  \footnote{\lr{ICJ. \textit{Accordance with International Law of the Unilateral Declaration of Independence in Respect of Kosovo}. Advisory Opinion, 22 July 2010. ICJ Reports 2010, p.~403.}}
  \item \textbf{کریمه (۲۰۱۴):} الحاق توسط روسیه با ادعای «رفراندوم»؛ محکوم‌شده توسط مجمع عمومی.%
  \footnote{\lr{UN General Assembly Resolution 68/262. 27 March 2014.}}
  \item \textbf{کاتالونیا (۲۰۱۷):} رفراندوم غیرقانونی از نظر دادگاه قانون اساسی اسپانیا.
  \item \textbf{اقلیم کردستان عراق (۲۰۱۷):} رفراندوم استقلال؛ عملاً بی‌نتیجه.
\end{itemize}

% ─────────────────────────────────────────────────────────────────────────
\section{مداخله‌ی انسان‌دوستانه: مرز باریک حقوق و اخلاق}
\label{sec:ch02-humanitarian}
% ─────────────────────────────────────────────────────────────────────────

\subsection{تعریف و ویژگی‌ها}
\label{subsec:ch02-hum-def}

\concept{مداخله‌ی انسان‌دوستانه} (\lr{Humanitarian Intervention}) به «توسل به زور توسط یک دولت، گروهی از دولت‌ها، یا سازمان بین‌المللی، عمدتاً به‌منظور حمایت از حقوق بشر اتباع دولت هدف» تعریف می‌شود.%
\footnote{\lr{Welsh, Jennifer M. \textit{Humanitarian Intervention and International Relations}. Oxford UP, 2004, p.~3.}}

\begin{tcolorbox}[legalbox, title={معیارهای مشروعیت مداخله‌ی انسان‌دوستانه (بر اساس گزارش \lr{ICISS})}]
\begin{enumerate}[nosep, label=\textbf{\arabic*.}]
  \item \textbf{آستانه‌ی عادلانه} (\lr{Just Cause}): نسل‌کشی، پاک‌سازی قومی، یا تلفات انسانی گسترده.
  \item \textbf{نیت درست} (\lr{Right Intention}): هدف اصلی باید توقف رنج بشری باشد.
  \item \textbf{آخرین راه‌حل} (\lr{Last Resort}): تمامی ابزارهای مسالمت‌آمیز امتحان شده باشند.
  \item \textbf{ابزار متناسب} (\lr{Proportional Means}): اقدام نظامی متناسب با تهدید باشد.
  \item \textbf{چشم‌انداز معقول} (\lr{Reasonable Prospects}): احتمال معقولی وجود داشته باشد که مداخله وضعیت را بهبود بخشد نه بدتر کند.
  \item \textbf{مرجع صالح} (\lr{Right Authority}): ترجیحاً شورای امنیت سازمان ملل.
\end{enumerate}
\vspace{4pt}
\textit{منبع:} \lr{ICISS, 2001, pp.~32--37.}
\end{tcolorbox}

\subsection{نقد اخلاقی-سیاسی معیارهای مشروعیت}
\label{subsec:ch02-hum-critique}

\thinker{مایکل والزر} (\lr{Michael Walzer}) در اثر کلاسیک «جنگ‌های عادلانه و ناعادلانه» به تنش بنیادین میان حاکمیت و حقوق بشر پرداخته و استدلال می‌کند که مداخله‌ی نظامی تنها در سه مورد مجاز است: جدایی‌طلبی مشروع، ضدحمله‌ی متقابل، و توقف نسل‌کشی یا فاجعه‌ی بشری.%
\footnote{\lr{Walzer, Michael. \textit{Just and Unjust Wars}. 5th ed. Basic Books, 2015, pp.~86--108.}}

در مقابل، \thinker{آلن بوکانان} (\lr{Allen Buchanan}) از حق مداخله‌ی «گسترده‌تر» دفاع می‌کند و استدلال می‌کند که اصل حاکمیت تنها زمانی معتبر است که دولت حداقلی از عدالت را رعایت کند.%
\footnote{\lr{Buchanan, Allen. \textit{Justice, Legitimacy, and Self-Determination}. Oxford UP, 2004, pp.~331--400.}}

\begin{figure}[htbp]
\centering
\begin{tikzpicture}[
  node distance=0.6cm,
  every node/.style={font=\small},
  critbox/.style={
    draw=AccentRed!70, fill=BgFail!50, rounded corners=3pt,
    text width=6.5cm, align=right, inner sep=6pt, font=\small
  },
  defbox/.style={
    draw=AccentGreen!70, fill=BgSuccess!50, rounded corners=3pt,
    text width=6.5cm, align=right, inner sep=6pt, font=\small
  },
]

\node[font=\large\bfseries, color=AccentRed] (ct) {\rl{نقدها}};
\node[critbox, below=0.4cm of ct] (c1)
  {\rl{گزینشی بودن: چرا لیبی آری و سوریه نه؟}};
\node[critbox, below=0.3cm of c1] (c2)
  {\rl{ابزار تغییر رژیم به نفع قدرت‌های بزرگ}};
\node[critbox, below=0.3cm of c2] (c3)
  {\rl{نادیده‌گرفتن ریشه‌های ساختاری بحران}};
\node[critbox, below=0.3cm of c3] (c4)
  {\rl{نبود سازوکار پاسخ‌گویی پس از مداخله}};
\node[critbox, below=0.3cm of c4] (c5)
  {\rl{تشدید بی‌ثباتی بلندمدت}};

\node[font=\large\bfseries, color=AccentGreen, right=5.5cm of ct] (dt) {\rl{دفاعیه‌ها}};
\node[defbox, below=0.4cm of dt] (d1)
  {\rl{نجات جان انسان‌ها در بحران فوری}};
\node[defbox, below=0.3cm of d1] (d2)
  {\rl{تقویت هنجار مسئولیت دولت‌ها در قبال شهروندان}};
\node[defbox, below=0.3cm of d2] (d3)
  {\rl{ایجاد فشار بازدارنده بر دیکتاتورها}};
\node[defbox, below=0.3cm of d3] (d4)
  {\rl{وجود سازوکار نهادی (شورای امنیت)}};
\node[defbox, below=0.3cm of d4] (d5)
  {\rl{حق اخلاقی مردم تحت ستم برای طلب کمک}};

\end{tikzpicture}
\caption{نقدها و دفاعیه‌های مداخله‌ی انسان‌دوستانه}
\label{fig:ch02-hum-debate}
\end{figure}

% ─────────────────────────────────────────────────────────────────────────
\section{چارچوب حقوقی استمداد از قدرت خارجی برای رفع استبداد}
\label{sec:ch02-tyranny}
% ─────────────────────────────────────────────────────────────────────────

\subsection{سنت حقوق طبیعی و حق مقاومت}
\label{subsec:ch02-resistance}

پرسش از مشروعیت استمداد خارجی برای رفع استبداد، ریشه در سنت کهن \concept{حق مقاومت در برابر ظلم} (\lr{Right of Resistance}) دارد. این حق در آثار متفکران مختلف بازتاب یافته:

\begin{longtable}{
  >{\raggedright\arraybackslash\bfseries}p{2.5cm}
  >{\raggedright\arraybackslash}p{2.5cm}
  >{\raggedright\arraybackslash}p{8cm}}
\caption{سیر تاریخی حق مقاومت در برابر ظلم}\label{tab:ch02-resistance}\\
\toprule
\rowcolor{PrimaryDeep}
\textcolor{white}{\textbf{متفکر/سنت}} &
\textcolor{white}{\textbf{دوره}} &
\textcolor{white}{\textbf{موضع اصلی}} \\
\midrule
\endfirsthead
\toprule
\rowcolor{PrimaryDeep}
\textcolor{white}{\textbf{متفکر/سنت}} &
\textcolor{white}{\textbf{دوره}} &
\textcolor{white}{\textbf{موضع اصلی}} \\
\midrule
\endhead
\bottomrule
\endlastfoot

توماس آکویناس &
سده‌ی ۱۳ &
مقاومت در برابر حاکم ستم‌گر مجاز است مشروط بر آنکه ستم آشکار و بی‌درمان باشد%
\footnote{\lr{Aquinas, Thomas. \textit{Summa Theologica}. II-II, Q.~42, Art.~2.}} \\
\rowcolor{NeutralLight}
جان لاک &
سده‌ی ۱۷ &
اگر حاکم «قرارداد اجتماعی» را نقض کند، مردم حق شورش و حتی دعوت از یاری‌رسان خارجی دارند%
\footnote{\lr{Locke, John. \textit{Two Treatises of Government}. Ed. Peter Laslett. Cambridge UP, 1988 [1689], Second Treatise, Ch.~19, sec.~222.}} \\
فقه اسلامی &
متقدم و متأخر &
اختلاف‌نظر شدید: برخی خروج بر حاکم جائر را واجب و برخی آن را حرام می‌دانند. استمداد از «کافر» اجماعاً ممنوع%
\footnote{طباطبایی، سیدجواد. \textit{زوال اندیشه‌ی سیاسی در ایران}. کویر، ۱۳۸۳، ص ۲۱۵--۲۴۰.} \\
\rowcolor{NeutralLight}
فلسفه‌ی روشنگری &
سده‌ی ۱۸ &
حق مقاومت جزء حقوق طبیعی بشر (اعلامیه‌ی حقوق بشر و شهروند فرانسه، ماده‌ی ۲) \\
حقوق بین‌الملل مدرن &
قرن ۲۰--۲۱ &
حق مقاومت مسلحانه‌ی ملت‌های تحت استعمار (قطعنامه ۲۶۲۵) اما نه ملت‌های در دولت‌های مستقل%
\footnote{\lr{Cassese, Antonio. \textit{Self-Determination of Peoples}. Cambridge UP, 1995, pp.~150--170.}} \\

\end{longtable}

\subsection{معضل «در شهر هر آنکه هست گیرند»}
\label{subsec:ch02-pandora}

نکته‌ی بنیادین حقوقی آن است که \concept{مشروعیت‌بخشی به مداخله‌ی خارجی یک فرایند برگشت‌ناپذیر است}. هنگامی که دری به مداخله گشوده شود، کنترل حدود و ثغور آن از دست دعوت‌کنندگان خارج می‌شود.

\begin{tcolorbox}[enemybox, title={«گر رسم شود که مست گیرند، در شهر باید هر آنکه هست گیرند» — تحلیل حقوقی}]
این ضرب‌المثل فارسی، بیان‌گر اصل حقوقی مهمی است: \textbf{مداخله‌ی خارجی، حتی اگر با دعوت آغاز شود، ذاتاً تمایل به گسترش فراتر از محدوده‌ی اولیه دارد.} دلایل حقوقی این گسترش عبارت‌اند از:

\begin{enumerate}[nosep, label=\textbf{\arabic*.}]
  \item \textbf{ابهام در حدود رضایت:} رضایت اولیه به‌ندرت شامل تعریف دقیق حدود عملیات، مدت زمان، و شرایط خروج می‌شود.
  \item \textbf{پویایی جنگ:} عملیات نظامی منطق خاص خود را دارد و نمی‌توان آن را با دقت حقوقی هدایت کرد (\lr{fog of war}).%
  \footnote{\lr{Clausewitz, Carl von. \textit{On War}. Trans. Howard and Paret. Princeton UP, 1976, Book I, Ch.~7.}}
  \item \textbf{عدم تقارن قدرت:} مداخله‌گر معمولاً قوی‌تر از دعوت‌کننده است و در عمل، شرایط را تعیین می‌کند.
  \item \textbf{منافع مداخله‌گر:} مداخله‌گر ناگزیر منافع خود را نیز دنبال می‌کند — بنابراین حدود مداخله تابع منافع اوست نه خواسته‌ی دعوت‌کننده.
\end{enumerate}
\end{tcolorbox}

\thinker{جان استوارت میل} (\lr{John Stuart Mill}) در مقاله‌ی معروف «چند کلمه درباره‌ی عدم مداخله» (۱۸۵۹) استدلال می‌کند که حتی مداخله‌ی خیرخواهانه مضر است زیرا ملتی که آزادی‌اش را با نیروی بیگانه به‌دست آورد، توانایی حفظ آن را نخواهد داشت:

\begin{histQuote}{جان استوارت میل، ۱۸۵۹}
«ملتی که آزادی خود را مدیون دست بیگانه باشد ... معمولاً شایسته‌ی حفظ آن نیست. اگر ملتی نتواند آزادی‌اش را خود به‌دست آورد ... به‌ندرت می‌تواند آن را حفظ کند ... آزادی‌ای که از بیرون تحمیل شود، آزادی نیست بلکه شکلی دیگر از سلطه است.»%
\footnote{\lr{Mill, John Stuart. "A Few Words on Non-Intervention." \textit{Fraser's Magazine} 60 (December 1859): 766--776. Reprinted in \textit{Dissertations and Discussions}, vol.~3. London, 1867, pp.~153--178.}}
\end{histQuote}

% ─────────────────────────────────────────────────────────────────────────
\section{درخواست جدایی و حقوق بین‌الملل: تحلیل انتقادی}
\label{sec:ch02-secession}
% ─────────────────────────────────────────────────────────────────────────

\subsection{حق جدایی در حقوق بین‌الملل: موجود یا معدوم؟}
\label{subsec:ch02-secession-right}

بر خلاف تصور عمومی، حقوق بین‌الملل \textbf{حق عامی برای جدایی یک‌جانبه} (\lr{unilateral secession}) قائل نیست. دیوان عالی کانادا در نظریه‌ی مشورتی مشهور خود درباره‌ی جدایی کبک (۱۹۹۸) این نکته را به‌روشنی بیان کرد:

\begin{legalNote}{نظر مشورتی دیوان عالی کانادا درباره‌ی جدایی کبک (۱۹۹۸)}
«نه حقوق بین‌الملل و نه حقوق داخلی کانادا، حق جدایی یک‌جانبه‌ی کبک را به رسمیت نمی‌شناسند ... با این حال، حق تعیین سرنوشت ممکن است در شرایط استثنایی — یعنی هنگامی که مردمی از دسترسی معنادار به حکومت برای پیگیری توسعه‌ی سیاسی، اقتصادی، اجتماعی و فرهنگی خود محروم شوند — حق جدایی اصلاحی (\lr{remedial secession}) را ایجاد کند.»%
\footnote{\lr{Supreme Court of Canada. \textit{Reference Re Secession of Quebec}. [1998] 2 SCR 217, paras.~122--123, 134--135.}}
\end{legalNote}

\subsection{جدایی اصلاحی: نظریه‌ای در حال شکل‌گیری}
\label{subsec:ch02-remedial}

نظریه‌ی \concept{جدایی اصلاحی} (\lr{Remedial Secession}) بر آن است که اگر دولتی به‌طور سیستماتیک حقوق بنیادین بخشی از شهروندان خود را نقض کند و هیچ راه‌حل داخلی وجود نداشته باشد، آن گروه حق جدایی دارد. اما این نظریه هنوز در حقوق بین‌الملل عرفی تثبیت نشده است.%
\footnote{\lr{Vidmar, Jure. "Remedial Secession in International Law: Theory and (Lack of) Practice." \textit{St Antony's International Review} 6, no.~1 (2010): 37--56.}}

\begin{table}[htbp]
\centering
\caption{مقایسه‌ی نظریه‌های حقوقی درباره‌ی جدایی}\label{tab:ch02-secession-theories}
\begin{adjustbox}{max width=\linewidth}
\begin{tabular}{
  >{\raggedleft\arraybackslash\bfseries}p{3cm}
  >{\raggedleft\arraybackslash}p{3.5cm}
  >{\raggedleft\arraybackslash}p{3.5cm}
  >{\raggedleft\arraybackslash}p{3.5cm}}
\toprule
\rowcolor{PrimaryDeep}
\textcolor{white}{نظریه} &
\textcolor{white}{شرط جدایی} &
\textcolor{white}{طرفداران اصلی} &
\textcolor{white}{انتقاد اصلی} \\
\midrule
حق اولیه (\lr{Primary Right}) &
رأی اکثریت ساکنان منطقه کافی است &
\lr{Beran, Philpott} &
بی‌ثبات‌سازی نظم بین‌الملل \\
\rowcolor{NeutralLight}
جدایی اصلاحی (\lr{Remedial Only}) &
فقط در صورت نقض فاحش حقوق &
\lr{Buchanan, Cassese} &
ابهام در تعریف «نقض فاحش» \\
ممنوعیت مطلق &
جدایی هرگز مجاز نیست &
بسیاری از دولت‌ها (عملاً) &
نادیده‌گرفتن حق تعیین سرنوشت \\
\bottomrule
\end{tabular}
\end{adjustbox}
\end{table}

% ─────────────────────────────────────────────────────────────────────────
\section{نمودار تحلیلی: مسیرهای حقوقی مداخله}
\label{sec:ch02-flowchart}
% ─────────────────────────────────────────────────────────────────────────

\begin{figure}[htbp]
\centering
\begin{tikzpicture}[
  node distance=1.2cm and 1.5cm,
  every node/.style={font=\small},
  decision/.style={
    diamond, draw=PrimaryMid, fill=BgCool,
    text width=2.8cm, align=center,
    inner sep=2pt, aspect=2
  },
  process/.style={
    rectangle, draw=PrimaryLight, fill=BgCool!50,
    rounded corners=4pt, text width=3.2cm, align=center,
    minimum height=1cm
  },
  result/.style={
    rectangle, draw=AccentGreen, fill=BgSuccess!50,
    rounded corners=4pt, text width=3cm, align=center,
    minimum height=1cm
  },
  reject/.style={
    rectangle, draw=AccentRed, fill=BgFail!50,
    rounded corners=4pt, text width=3cm, align=center,
    minimum height=1cm
  },
  arrow/.style={->, >=Stealth, thick, PrimaryMid},
  ymark/.style={font=\scriptsize\bfseries, color=AccentGreen},
  nmark/.style={font=\scriptsize\bfseries, color=AccentRed},
]

\node[process] (start) {\rl{درخواست مداخله‌ی خارجی}};

\node[decision, below=1.2cm of start] (d1) {\rl{آیا حکومت مشروع%
 دعوت کرده؟}};
\draw[arrow] (start) -- (d1);

\node[process, left=2.5cm of d1] (inv) {\rl{مداخله بر اساس%
 دعوت}};
\draw[arrow] (d1) -- node[ymark, above] {\rl{بله}} (inv);

\node[decision, below=1.5cm of d1] (d2) {\rl{آیا نسل‌کشی/%
 جنایت علیه بشریت رخ داده؟}};
\draw[arrow] (d1) -- node[nmark, right] {\rl{خیر}} (d2);

\node[decision, left=2.5cm of d2] (d3) {\rl{آیا شورای%
 امنیت مجوز داده؟}};
\draw[arrow] (d2) -- node[ymark, above] {\rl{بله}} (d3);

\node[result, below=1.2cm of d3] (r2p) {\rl{مداخله بر اساس%
 \lr{R2P}}};
\draw[arrow] (d3) -- node[ymark, right] {\rl{بله}} (r2p);

\node[decision, below=1.5cm of d2] (d4) {\rl{آیا ملت تحت%
 استعمار/اشغال است؟}};
\draw[arrow] (d2) -- node[nmark, right] {\rl{خیر}} (d4);

\node[result, left=2.5cm of d4] (lib) {\rl{جنبش آزادی‌بخش%
 (حق تعیین سرنوشت)}};
\draw[arrow] (d4) -- node[ymark, above] {\rl{بله}} (lib);

\node[reject, below=1.2cm of d4] (rej) {\rl{مداخله غیرمجاز%
 در حقوق بین‌الملل}};
\draw[arrow] (d4) -- node[nmark, right] {\rl{خیر}} (rej);

\node[reject, below=1.2cm of d3, xshift=3cm] (grey) {\rl{منطقه‌ی خاکستری:%
 مداخله بدون مجوز}};
\draw[arrow] (d3) -- node[nmark, below, sloped] {\rl{خیر}} (grey);

\end{tikzpicture}
\caption{نمودار تصمیم‌گیری حقوقی مداخله‌ی خارجی}
\label{fig:ch02-flowchart}
\end{figure}

% ─────────────────────────────────────────────────────────────────────────
\section{جمع‌بندی فصل}
\label{sec:ch02-summary}
% ─────────────────────────────────────────────────────────────────────────

این فصل نشان داد که نظام حقوقی بین‌المللی در مورد مداخله‌ی خارجی، در تنش دائمی میان دو اصل بنیادین قرار دارد: \concept{حاکمیت دولت‌ها} و \concept{حقوق بشر}. از صلح وستفالی (۱۶۴۸) تا دکترین مسئولیت حمایت (۲۰۰۵)، هنجارهای بین‌المللی به‌تدریج فضای بیشتری برای مداخله باز کرده‌اند، اما این گشایش همواره با خطر سوءاستفاده همراه بوده است.

نکات کلیدی:
\begin{itemize}[nosep, rightmargin=1em]
  \item مداخله بر اساس دعوت حکومت مشروع، از نظر فنی-حقوقی مجاز اما اخلاقاً بحث‌برانگیز است.
  \item حق تعیین سرنوشت شامل حق جدایی یک‌جانبه نمی‌شود مگر در شرایط استثنایی (جدایی اصلاحی).
  \item دکترین \lr{R2P} محدود به چهار جرم بین‌المللی است و اِعمال آن منوط به تأیید شورای امنیت است.
  \item هشدار اساسی میل و ضرب‌المثل «مست گیرند...» حقیقتی تجربی است: مداخله ذاتاً گسترش‌یابنده است.
\end{itemize}

\cleardoublepage